\chapter{The statistics of the experiments}
\label{ch:experimental-statistics}

Every number in the paper's tables and figures is a statistic of a small number of
replicates, sixteen at the frozen configuration and eight in the sweeps, and every
comparison between two arms is paired on identical data. This chapter defines the
quantities the paper reports (RMSE, bias, dispersion, standard errors, bootstrap
intervals, paired differences and ratios, log-slopes), explains why a local slope
over a finite window is not an asymptotic exponent, describes the bundle structure
(training, validation, test), the two kinds of seeds (development and confirmatory),
and what ``pre-registered'' and ``pre-data amendment'' mean in this project. The
numbers themselves are in the paper; here are the definitions of what they are
numbers \emph{of}.

\section{Replicates, seeds and the inferential unit}

A \term{replicate} of an experiment is one complete run with one value of the random
\term{seed}: the seed determines the clean sequences of the training, validation and
test bundles, the noise draws that corrupt them, the initialisation of the estimator,
and any Monte Carlo inside the fit. Two arms run at the same seed see the same data,
which is what makes their comparison paired (\cref{sec:paired}). The \term{inferential
unit} is the replicate: every standard error, interval and sign count below is over
seeds, never over sequences or sites within a seed, because sequences within a seed
share the bundle and sites within a sequence are dependent.

Let $e_1,\dots,e_n$ be the values of a statistic (an error, a fitted coordinate) over
$n$ seeds. The \term{mean over seeds} is $\bar e=\tfrac1n\sum_se_s$, the \term{sample
standard deviation} is $\hat\sigma=\bigl(\tfrac1{n-1}\sum_s(e_s-\bar e)^2\bigr)^{1/2}$,
and the \term{standard error of the mean} is $\hat\sigma/\sqrt n$. With $n=16$ the
standard error is a quarter of the standard deviation; with $n=8$ it is $0.35$ of it.

\section{RMSE, bias and dispersion}

Let $\hat\theta_1,\dots,\hat\theta_n$ be estimates of a known truth $\theta$ over $n$
seeds (the paper's fitted correlation and fitted excess kurtosis at each level of
E06 are of this kind). The \term{root mean squared error} over replicates is
\begin{equation}
  \mathrm{RMSE}=\Bigl(\frac1n\sum_{s=1}^n(\hat\theta_s-\theta)^2\Bigr)^{1/2},
\end{equation}
and it decomposes into a \term{bias} and a \term{dispersion}:
\begin{equation}
  \frac1n\sum_s(\hat\theta_s-\theta)^2=\underbrace{(\bar\theta-\theta)^2}_{\text{bias}^2}
  +\underbrace{\frac1n\sum_s(\hat\theta_s-\bar\theta)^2}_{\text{dispersion}^2},
  \qquad \bar\theta=\frac1n\sum_s\hat\theta_s ,
  \label{eq:rmse-decomp}
\end{equation}
because the cross term $\tfrac2n\sum_s(\hat\theta_s-\bar\theta)(\bar\theta-\theta)$
vanishes. The dispersion is the population-form standard deviation (divisor $n$);
the sample form (divisor $n-1$) differs by the factor $\sqrt{n/(n-1)}$, which is
$1.03$ at $n=16$. The paper reports RMSE, bias and dispersion separately at each level
because they answer different questions: the Cram\'er--Rao reference of the chapter on
Fisher information bounds the dispersion of an unbiased estimator, so an estimator
whose bias is not zero can have a dispersion below the reference, and only the
decomposition says which is happening.

\paragraph{The sampling error of a standard deviation.}
A dispersion estimated from $n$ replicates is itself random. If the replicates were
normal, $(n-1)\hat\sigma^2/\sigma^2$ would be $\chi^2_{n-1}$, whose relative standard
deviation is $\sqrt{2/(n-1)}$, so the relative standard error of $\hat\sigma$ is about
$1/\sqrt{2(n-1)}$: $1/\sqrt{30}\approx18\%$ at $n=16$ and $1/\sqrt{14}\approx27\%$ at
$n=8$. This is what the paper means by ``within the sampling error of a
sixteen-replicate standard deviation'': two dispersions that differ by less than about
a fifth are not distinguishable at this replicate count.

\section{Confidence intervals by the bootstrap}
\label{sec:bootstrap}

The paper's intervals are \term{percentile bootstrap} intervals over seeds. The
principle is the plug-in principle: the unknown law of the statistic under
resampling of the true population is approximated by its law under resampling of the
observed replicates.

\begin{definition}[Percentile bootstrap over seeds]
Given per-seed values $e_1,\dots,e_n$ and a statistic $T(e_1,\dots,e_n)$ (a mean, a
mean of paired differences, a slope), draw $B$ \term{resamples}
$(e_{s_1^*},\dots,e_{s_n^*})$ with the indices $s_j^*$ drawn independently and uniformly
from $\{1,\dots,n\}$ with replacement, compute $T^*_b$ for each, and report the
interval $[T^*_{(0.025)},T^*_{(0.975)}]$ of the $2.5$th and $97.5$th percentiles of
$\{T^*_b\}_{b=1}^B$.
\end{definition}

The code uses $B=20{,}000$ resamples with a fixed bootstrap seed; the function is
\code{boot\_ci} in \code{tools/make\_capacity\_numbers.py}, and the other number
generators carry copies of it. Two things must be said
about such an interval. It carries \emph{Monte Carlo error} from the finite $B$, which
is largest in the tails: the $97.5$th percentile of $20{,}000$ draws rides on the
$500$ largest values, and with $n=8$ seeds the bootstrap distribution is discrete with
large atoms, so the percentile can jump between neighbouring atoms from one stream
to another. And it is not an exact confidence interval: with $n=8$ or $16$ the
percentile interval tends to under-cover, and it inherits any skewness of the
statistic. The paper therefore reports the interval as the outcome of a declared
procedure and reads it as such, never as a $p$-value.

\paragraph{The exact enumerated bootstrap.}
A resample is a multiset of $n$ draws from $n$ seeds, and the statistic depends only
on the multiset (the counts $c_1,\dots,c_n$ with $\sum_sc_s=n$) when it is symmetric in
the replicates. The number of distinct multisets is the number of weak compositions
of $n$ into $n$ parts,
\begin{equation}
  \binom{2n-1}{n},\qquad\text{which is }6435\text{ at }n=8,\quad 92{,}378\text{ at }n=10,\quad 300{,}540{,}195\text{ at }n=16,
\end{equation}
and the multiset with counts $(c_s)$ has probability $\dfrac{n!}{\prod_sc_s!}\,n^{-n}$
under resampling with replacement (the multinomial law). When the count is small
enough, the bootstrap distribution is therefore \emph{enumerable}: compute $T$ on every
multiset, weight it by its multinomial probability, and read the weighted percentiles.
The interval then has no Monte Carlo error at all. The tool
\code{tools/make\_sample\_complexity.py} does this whenever $\binom{2n-1}{n}\le200{,}000$
(that is, $n\le10$) and falls back to Monte Carlo with a reported $B$ above that limit.
This is the \term{exact enumerated bootstrap} the paper's provenance appendix refers to,
and it entered the project because a Monte Carlo interval for a fitted slope was found
to have its upper limit on one side or the other of the value being tested depending
on the bootstrap stream (\code{docs/SCIENCE\_AUDIT.md}, finding 3 of 18 September);
enumeration settled it. At $n=16$ the enumeration is out of reach and $B=20{,}000$
Monte Carlo draws are used, which is adequate for a mean over sixteen seeds.

\section{Paired differences and paired ratios}
\label{sec:paired}

Two arms $A$ and $B$ run at the same seed are fitted on the same training bundle,
selected on the same validation bundle and scored on the same test bundle. Their
errors $e^A_s$ and $e^B_s$ are then \term{paired}, and the comparison is made on the
per-seed differences or ratios, not on the two means.

\paragraph{Paired differences.}
With $d_s=e^A_s-e^B_s$, the mean paired difference is $\bar d=\bar e^A-\bar e^B$, the
same number as the difference of means; what changes is its uncertainty:
\begin{equation}
  \Var(\bar d)=\frac1n\bigl(\Var(e^A)+\Var(e^B)-2\Cov(e^A,e^B)\bigr),
\end{equation}
and the covariance is large and positive when both arms suffer from the same hard
seeds (a bundle with an unusual chain, a heavy-tailed draw). Pairing removes that
shared variation. The paper reports $\bar d$ with a percentile bootstrap of the $d_s$
over seeds, and beside it the \term{sign count}, ``in favour of $A$ in $k$ of $n$
seeds'', a distribution-free summary that needs no scale.

\paragraph{Paired ratios.}
When the question is a factor (``the best network is a factor $r$ behind''), the paper
forms the per-seed ratio $r_s=e^A_s/e^B_s$ and reports the \term{mean of ratios}
$\bar r=\tfrac1n\sum_sr_s$ with a percentile bootstrap. This is not the same as the
\term{ratio of means} $\bar e^A/\bar e^B$: for positive quantities the mean of ratios
exceeds the ratio of means when the denominator varies across seeds (Jensen's
inequality applied to $1/e^B$), and the two can differ by a few per cent at $n=16$.
The paper states which it uses (``ratios are formed within seed and averaged'') and,
for the headline factor, reports both, so that a reader can see the size of the
difference. A paired ratio is only meaningful when both arms are scored on the same
test bundle; the generator that joins rows from different experiments checks the
SHA-256 of the bundles and raises if they differ (\code{docs/EXPERIMENTAL\_CONTRACT.md},
Section 3).

\section{Log-slopes and why a local slope is not an asymptotic exponent}
\label{sec:slopes}

Many of the paper's theoretical statements are power laws in $e^{-t}$: an information
decays as $e^{-4t}$ or $e^{-8t}$, an RMSE grows as $e^{2t}$ or $e^{4t}$. A measured curve
$y(t)$ is compared with such a law through its \term{log-slope}, the derivative of
$\log y$ with respect to $t$, so that $y\propto e^{-8t}$ has log-slope $-8$.

\paragraph{The OLS log-slope over a window.}
Given values $y_1,\dots,y_K$ at levels $t_1<\dots<t_K$ in a declared window, the
ordinary least-squares slope of $\log y$ on $t$ is
\begin{equation}
  \hat b=\frac{\sum_k(t_k-\bar t)(\log y_k-\overline{\log y})}{\sum_k(t_k-\bar t)^2},
  \label{eq:ols-slope}
\end{equation}
which for $K=2$ reduces to the two-point slope $(\log y_2-\log y_1)/(t_2-t_1)$. The
window and the estimator are declared in the registry entry before the data are seen
(\cref{sec:prereg}); the two-point and the OLS estimators differ, and the audit of 18
September replaced a two-point estimate by the pre-declared OLS one for exactly this
reason. Its uncertainty is a bootstrap over seeds of the whole fit
(\cref{sec:bootstrap}), because the levels within a seed share the seed's chains.

\paragraph{Local versus asymptotic.}
An asymptotic law $y(t)=Ce^{-\beta t}(1+\epsilon_t)$ with $\epsilon_t\to0$ has log-slope
\begin{equation}
  \frac{\dd\log y}{\dd t}=-\beta+\frac{\dd}{\dd t}\log(1+\epsilon_t),
\end{equation}
and the second term is not small until $\epsilon_t$ is. For the paper's Gaussian chain
$\epsilon_t=O(e^{-2t}\lmax(\Sig_0))$, and the scalar example of the chapter on Fisher
information makes the effect explicit: for a noised Gaussian coordinate of variance
$\theta$, $I_t=\tfrac12e^{-4t}(1+(\theta-1)e^{-2t})^{-2}$ has local log-slope
\begin{equation}
  -4+\frac{4(\theta-1)e^{-2t}}{1+(\theta-1)e^{-2t}},
\end{equation}
which for $\theta=5$ is about $-1.6$ at $t=0.5$, $-3.3$ at $t=1.5$ and $-4.0$ at $t=3$.
A window that ends near the crossover $\tfrac12\log(\theta-1)\approx0.7$ measures a
slope of $-2$ or so and says nothing against the asymptotic $-4$; a window well above it
measures $-4$. The OLS slope over a window is a weighted average of the local slopes
in it, so it lies between the shallowest and the steepest local values. This is the
whole content of the paper's remark that slopes of about $-3$ and $-6$ measured over
$t\le1.6$ at $\corr=0.85$ are pre-asymptotic values below $\tstar$, and of the
pre-data amendment of E12b, which replaced ``compare the fitted slope with $4$ and
$8$'' by ``compare $M^*(t)$ with the exact information curve, whose local slope on the
swept levels is what the corollary predicts there''.

\section{Bundles: training, validation and test}
\label{sec:bundles}

Each seed owns three \term{bundles} of sequences drawn from three disjoint random
streams keyed by a tag, the seed and, for training, the size $\nseq$
(\code{rng\_for} in \code{src/markov\_score/utils.py}; the tags are in
\code{experiments/protocol.py}):
\begin{itemize}
\item the \term{training bundle}, $\nseq\in\{32,128,512,2048\}$ sequences, the only
  data the estimator's fit sees; under the one-view protocol each sequence
  contributes one noised view at one level, the sequences being split evenly across
  the twelve levels;
\item the \term{validation bundle}, $1024$ sequences, used only to \emph{select}: the
  EM iteration count on its checkpoint ladder, the network checkpoint, and, for the
  architecture screens, nothing (those use development seeds, \cref{sec:seeds});
\item the \term{test bundle}, $2048$ sequences, used only to \emph{report}.
\end{itemize}
Selection on validation and reporting on test are separated so that the reported
number is not optimistically biased by the choice it depended on. The regions differ
too: selection uses the \term{bulk} (sites $5$ to $\nsites-4$), reporting the headline
uses all sites, and the centre site is a diagnostic. The SHA-256 of every bundle is
stored with every row, so that pairing across arms and across experiments is
checkable after the fact.

\paragraph{The estimand.}
The headline quantity is the \term{schedule-pooled relative score error},
\begin{equation}
  \Bigl(\frac{\sum_t\sum_{j,i}\bigl(\hat s_t(\vx_j)_i-s_t(\vx_j)_i\bigr)^2}{\sum_t\sum_{j,i}s_t(\vx_j)_i^2}\Bigr)^{1/2},
  \label{eq:estimand}
\end{equation}
the sum over the twelve levels $t$ of the schedule, the test sequences $j$ and the
sites $i$ of the squared error between an estimator's score $\hat s$ and the tree-exact
reference $s$, divided by the same sum of the reference's squared norm, and then
rooted. ``Pooled'' means that the sums are formed before the ratio and the root: the
low levels, where the score is large, dominate the pooled energy, and a per-level
statement (such as the Edgeworth correction of E16 or the LMMSE floor per level of
E07) must not be confused with a pooled one. A relative error of $x$ means the
estimator's squared error is $x^2$ of the reference's energy over the schedule.

\section{Development and confirmatory seeds}
\label{sec:seeds}

Two kinds of seed never mix. \term{Development seeds} ($10000$ to $10003$) are used for
every choice that a human or a script makes among alternatives: architectures,
learning rates, output parameterisations, the mixture capacity ladder's shape. Numbers
computed on them appear in no reported table. \term{Confirmatory seeds} ($0$ to $15$,
or $0$ to $7$ in the sweeps) are used for every reported number and for no choice.
The reason is selection bias: a configuration chosen because it did best on a set of
seeds has an error on those seeds that is biased downward (the winner's curse), and
only seeds that played no part in the choice give an unbiased reading. The
architecture screen of the paper's Appendix C is therefore a table on development
seeds whose winners are then run once on the confirmatory seeds. The per-seed
validation bundle is a third thing: it belongs to a confirmatory seed but is used only
to pick the checkpoint within an arm, which is part of the arm's definition rather
than a choice among arms.

\section{Pre-registration and pre-data amendments}
\label{sec:prereg}

A \term{registered experiment} is an entry of \code{experiments/registry/} (one JSON
file per identifier, E01 to E21) that records, before the experiment runs, the
hypothesis it tests with a stated prediction, the design (the one factor varied and
everything held fixed), the pre-declared analysis (which script computes which
statistic from which file, with which estimator and window), the seeds, and the
status. When the run is complete, its output files are \term{frozen}: their SHA-256
digests are written into the entry (\code{tools/freeze.py}), any later byte change
fails the build (\code{tools/freeze.py --check}, run by \code{tools/check\_manuscript.py}),
and a generator refuses to emit a number from an output that is not frozen or whose
provenance records a dirty code tree. Every number in the paper is a macro emitted by
such a generator from a frozen file, and an independent tool
(\code{tools/audit\_claims.py}) recomputes the headline numbers from the same frozen
files by a different code path. The contract is written out in
\code{docs/EXPERIMENTAL\_CONTRACT.md}.

A \term{pre-data amendment} is a change to a registered entry made \emph{before any
cell of the experiment has run}, recorded in the entry's \code{amendments} list with a
timestamp, what changed, why, and what did not change. It is the one kind of change
the contract allows to an entry; once a confirmatory cell has run, the entry is never
changed, and a fix is a new experiment with a new identifier or a versioned directory,
with the old run kept and reported for what it is. Two examples fix the idea.
\begin{itemize}
\item E14 (\code{experiments/registry/E14.json}) was pre-registered on 18 September
  and submitted to the cluster; while its jobs were still pending and no cell had run,
  the entry was amended so that each network row records the whole held-out objective
  curve rather than only its minimum, so that a diverged fit can be reported as a
  divergence and not as the network's error. The amendment states that nothing in the
  design, the arms, the seeds or the analysis changed, and the jobs were resubmitted.
\item E12b (\code{experiments/registry/E12b.json}) was amended, before any cell ran,
  to compare the measured sample-complexity curve with the exact information curve
  computed by Fisher's identity (E21) rather than with the asymptotic exponents $4$
  and $8$, because a computation made before the data showed that the local slope of
  the information curve on the swept levels is between $6$ and $7$ for the shape
  coordinate; the original plan would have recorded a confirmation as a refutation.
\end{itemize}
The value of the discipline is that a reader can tell, for every number, whether it
was predicted or observed, and whether the analysis that produced it was chosen before
or after the data existed. Where the paper reports something it did not pre-register
(the audit computations E15 to E21 were registered on the day they were run), it says so.

\section{Worked example: eight seeds}

Take $n=8$ confirmatory seeds, as in the regime sweeps. The exact bootstrap
enumerates $\binom{15}{8}=6435$ multisets. The multiset in which every seed appears
once has probability $8!/8^8=40320/16777216\approx0.0024$; a given seed is absent from
a resample with probability $(7/8)^8\approx0.344$, so a typical resample contains
about $8(1-0.344)\approx5.2$ distinct seeds, which is why bootstrap intervals at $n=8$
are wide and discrete. The relative standard error of an eight-replicate standard
deviation is about $1/\sqrt{14}\approx27\%$.

For a paired comparison, suppose four seeds gave errors $e^A=(0.10,0.12,0.11,0.13)$
and $e^B=(0.08,0.10,0.09,0.11)$. The unpaired means are $0.115$ and $0.095$, each with
a sample standard deviation of about $0.013$, and the unpaired standard error of their
difference is about $0.009$, so the two arms look barely distinguishable; the paired
differences are all exactly $0.02$, with zero dispersion, and $A$ is behind in $4$ of
$4$ seeds. The per-seed ratios are $1.25$, $1.20$, $1.22$ and $1.18$, with mean of ratios
$1.214$; the ratio of means is $0.115/0.095\approx1.211$, a little smaller, as
\cref{sec:paired} predicts. All of these are properties of the four invented numbers,
not of any experiment.

\where{{\sloppy The paper's Section 5 (paired seeds as the inferential unit, ratios formed
within seed, sixteen and eight seeds), Section 4 (RMSE, dispersion and bias per level,
the sixteen-replicate standard deviation, fitted log-slopes on declared windows),
Appendix C (bundles, seeds, the estimand, the checkpoint ladders, percentile bootstraps
over seeds, provenance) and Appendix E. Experiments: every one of E01 to E21; the
examples above are E06, E12, E12b, E14 and E19. Code: \code{experiments/protocol.py}
(sizes, seeds, bundle tags, the estimand and regions), \code{experiments/frozen\_config.py},
\code{src/markov\_score/utils.py} (\code{rng\_for}), \code{tools/make\_capacity\_numbers.py}
(\code{boot\_ci}), \code{tools/make\_sample\_complexity.py} (the enumerated bootstrap and
the OLS slope), \code{tools/make\_learning\_numbers.py}, \code{tools/freeze.py},
\code{tools/registry.py}, \code{tools/check\_manuscript.py}, \code{tools/audit\_claims.py},
\code{docs/EXPERIMENTAL\_CONTRACT.md}, \code{docs/SCIENCE\_AUDIT.md},
\code{experiments/registry/E14.json}, \code{experiments/registry/E12b.json}.\par}}

\paragraph{References.}
B. Efron and R. J. Tibshirani (1993), \emph{An Introduction to the Bootstrap}, Chapman
and Hall (Chapters 6 and 13: the bootstrap standard error and the percentile interval).
A. C. Davison and D. V. Hinkley (1997), \emph{Bootstrap Methods and Their Application},
Cambridge University Press (Chapter 2: the enumeration of resamples and the
multinomial weights; Chapter 5: percentile intervals and their coverage).
P. Hall (1992), \emph{The Bootstrap and Edgeworth Expansion}, Springer (Chapter 3: the
accuracy of percentile intervals).
L. Wasserman (2004), \emph{All of Statistics}, Springer (Chapters 6 to 9: bias,
variance, mean squared error, standard errors and the bootstrap).
B. A. Nosek, C. R. Ebersole, A. C. DeHaven and D. T. Mellor (2018), ``The
preregistration revolution'', Proceedings of the National Academy of Sciences 115,
2600--2606.
